\begin{proposition}
\label{lambda-pnt}
\lean{lambda_pnt}
% at /home/runner/work/trough/trough/library/PrimeNumberTheoremAnd/Consequences.lean:2201.0-2251.97
\leanok
We have $\sum_{n \leq x} \lambda(n) = o(x)$.
\end{proposition}
\begin{proof}
\uses{WeakPNT,mu-pnt,mun}
\leanok
From the identity
    $$ \lambda(n) = \sum_{d^2|n} \mu(n/d^2)$$
  and summing, we have
  $$ \sum_{n \leq x} \lambda(n) = \sum_{d \leq \sqrt{x}} \sum_{n \leq x/d^2} \mu(n).$$
  For any $\eps>0$, we have from Proposition \ref{mu-pnt} that
  $$ \sum_{n \leq x/d^2} \mu(n) = O(\eps x/d^2) + O_\eps(1)$$
  and hence on summing in $d$
  $$ \sum_{n \leq x} \lambda(n) = O(\eps x) + O_\eps(x^{1/2}).$$
  Sending $\eps \to 0$ we obtain the claim.
\end{proof}
